\PassOptionsToPackage{table,xcdraw}{xcolor}
\documentclass[aspectratio=169]{beamer}

\usetheme{Madrid}
\usecolortheme{default}
\definecolor{ucbblue}{RGB}{0, 51, 102}

\setbeamercolor{structure}{fg=ucbblue}
\setbeamercolor*{palette primary}{fg=white,bg=ucbblue}
\setbeamercolor*{palette secondary}{fg=white,bg=ucbblue!80}
\setbeamercolor*{palette tertiary}{fg=white,bg=ucbblue!90}
\setbeamercolor{title}{fg=white, bg=ucbblue}
\setbeamercolor{frametitle}{fg=white, bg=ucbblue}
\setbeamercolor{item}{fg=ucbblue}

\usepackage[T1]{fontenc}
\usepackage[utf8]{inputenc}
\usepackage{tgpagella}
\usefonttheme{serif}
\usepackage[spanish, es-noshorthands]{babel}
\usepackage{amsmath, amssymb, bm}
\usepackage{graphicx}
\usepackage[most]{tcolorbox}
\usepackage{tikz}
\usepackage[american]{circuitikz}
\usetikzlibrary{shapes.geometric, arrows.meta, positioning, calc}

\title[Cierre Hito 2]{Cierre Conceptual Hito 2: Modo Común, CMRR e Integración}
\subtitle{IMT-221 Circuitos Electrónicos III — Semana 9, Sesión 1}
\author[B. Quiroga Turdera]{M.Sc. Ing. Bernardo Quiroga Turdera}
\institute[UCB Tarija]{Universidad Católica Boliviana ``San Pablo'' — Sede Tarija}
\date{Semana 9}

\begin{document}

\begin{frame}
    \titlepage
\end{frame}

\begin{frame}{La Cadena de Razonamiento Integrada}
    \begin{tcolorbox}[colback=ucbblue!5, colframe=ucbblue, title=\textbf{El Núcleo Conceptual del Hito 2}]
        La evaluación no medirá memorización de fórmulas aisladas, sino la capacidad de leer datos, escoger el modelo, justificar supuestos y diagnosticar a través de toda la cadena.
    \end{tcolorbox}
    
    \vspace{0.5cm}
    \begin{center}
    \resizebox{0.95\textwidth}{!}{
    \begin{tikzpicture}[auto, thick]
        \node[draw, fill=blue!10, rounded corners, align=center] (dc) {Bias DC};
        \node[draw, fill=blue!10, rounded corners, align=center, right=0.4cm of dc] (q) {Punto Q};
        \node[draw, fill=green!10, rounded corners, align=center, right=0.4cm of q] (gm) {$g_m$};
        \node[draw, fill=green!10, rounded corners, align=center, right=0.4cm of gm] (ad) {$A_{d}$};
        \node[draw, fill=orange!10, rounded corners, align=center, right=0.4cm of ad] (vd) {$v_d / v_{cm}$};
        \node[draw, fill=orange!10, rounded corners, align=center, right=0.4cm of vd] (ac) {$A_{c}$};
        \node[draw, fill=red!10, rounded corners, align=center, right=0.4cm of ac] (cmrr) {CMRR};
        \node[draw, fill=red!20, rounded corners, align=center, right=0.4cm of cmrr] (int) {Interpretación};
        
        \draw[->] (dc) -- (q); \draw[->] (q) -- (gm); \draw[->] (gm) -- (ad);
        \draw[->] (ad) -- (vd); \draw[->] (vd) -- (ac); \draw[->] (ac) -- (cmrr); \draw[->] (cmrr) -- (int);
    \end{tikzpicture}
    }
    \end{center}
\end{frame}

\begin{frame}{Descomposición Diferencial y Modo Común}
    Cualquier par de entradas ($v_1, v_2$) puede separarse matemáticamente[cite: 7]:
    \begin{align*}
        \mathbf{v_d} &= \mathbf{v_1 - v_2} \quad \text{(Separación entre entradas)} \quad \text{[cite: 7]} \\
        \mathbf{v_{cm}} &= \mathbf{\frac{v_1 + v_2}{2}} \quad \text{(Nivel compartido por ambas)} \quad \text{[cite: 7]}
    \end{align*}
    
    \vspace{0.1cm}
    \begin{columns}[c]
        \begin{column}{0.48\textwidth}
            \textbf{Ejemplo: Entrada Mixta[cite: 7]} \\
            $v_1 = 106\,\text{mV}, \quad v_2 = 94\,\text{mV}$ \\
            $\implies v_d = 12\,\text{mV}$ \\
            $\implies v_{cm} = 100\,\text{mV}$
        \end{column}
        \begin{column}{0.48\textwidth}
            \begin{tcolorbox}[colback=yellow!10, colframe=orange, title=Nota Conceptual]
                $v_1 = v_2 \neq 0$ \textbf{no} significa ausencia de señal, significa Modo Común puro[cite: 7]. El modo común no es automáticamente "ruido"[cite: 7].
            \end{tcolorbox}
        \end{column}
    \end{columns}
\end{frame}

\begin{frame}{Comportamiento Físico bajo Modo Común ($v_1 = v_2 = v_{cm}$)}
    \begin{columns}[c]
        \begin{column}{0.4\textwidth}
            \centering
            \begin{circuitikz}[scale=0.75, transform shape, thick]
                \draw (0,0) node[npn](Q1){}; \node[left] at (Q1.B) {$v_{cm}$};
                \draw (3,0) node[npn, xscale=-1](Q2){}; \node[right] at (Q2.B) {$v_{cm}$};
                \draw (Q1.C) to[R, l=$R_C$] (0,2) -- (3,2) to[R, l_=$R_C$] (Q2.C);
                \draw (1.5,2) -- (1.5,2.5) node[above]{$+V_{CC}$};
                \draw (Q1.E) -- (0,-1) -- (3,-1) -- (Q2.E);
                \draw (1.5,-1) to[R, l_=$R_T$] (1.5,-3) node[below]{$-V_{EE}$};
                \draw[->, red] (1.7, -1.2) -- (1.7, -2.8) node[midway, right]{$I_T$};
            \end{circuitikz}
        \end{column}
        \begin{column}{0.55\textwidth}
            \textbf{¿Por qué hay ganancia $A_c$?[cite: 7]}
            \begin{itemize}
                \item \textbf{Fuente Ideal ($R_T \to \infty$):} $I_T$ es constante. Las corrientes de rama no cambian. Rechazo perfecto[cite: 7].
                \item \textbf{Fuente Real ($R_T < \infty$):} El movimiento de los emisores varía la caída sobre $R_T$[cite: 7]. \\
                La corriente total fluctúa: $\Delta I_T \neq 0$[cite: 7].
            \end{itemize}
            
            \vspace{0.2cm}
            Esta fluctuación se divide simétricamente:
            \begin{equation*}
                \Delta i_{C1} \approx \Delta i_{C2} \implies \Delta v_{C1} \approx \Delta v_{C2} \quad \text{[cite: 7]}
            \end{equation*}
            \textit{Ambos colectores se mueven juntos. La salida Single-Ended observa este error[cite: 7].}
        \end{column}
    \end{columns}
\end{frame}

\begin{frame}{Ganancia de Modo Común Single-Ended ($A_{c,se}$)}
    Para la convención oficial $\mathbf{v_o = v_{C2}}$[cite: 7]:
    
    Analizamos un \textbf{medio circuito} de modo común con resistencia de emisor equivalente a $2R_T$ (dado que la cola soporta $2i_e$)[cite: 7].
    Por LVK en la malla de entrada (Modelo T):
    \begin{equation*}
        v_{cm} = i_e r_e + 2 i_e R_T = i_e (r_e + 2R_T) \implies \mathbf{i_e = \frac{v_{cm}}{r_e + 2R_T}} \quad \text{[cite: 7]}
    \end{equation*}
    
    Como $v_o = -i_c R_C$ e $i_c \approx \alpha i_e$[cite: 7]:
    \begin{equation*}
        \mathbf{A_{c,se} = \frac{v_{C2}}{v_{cm}} \approx -\frac{\alpha R_C}{r_e + 2R_T}} \quad \text{[cite: 7]}
    \end{equation*}
    
    Si $R_T \gg r_e$ y $\alpha \approx 1$:
    \begin{equation*}
        \mathbf{A_{c,se} \approx -\frac{R_C}{2R_T}} \implies \text{A mayor } R_T \text{, menor } |A_c| \text{ y mejor CMRR[cite: 7].}
    \end{equation*}
\end{frame}

\begin{frame}{CMRR y la Puerta de Consistencia (Validity Gate)}
    \begin{alertblock}{Definición Oficial Hito 2}
        \begin{center}
            $\displaystyle \mathbf{CMRR = \left|\frac{A_{d,se}}{A_{c,se}}\right|} \qquad \mathbf{CMRR_{dB} = 20\log_{10}(CMRR)}$ \quad \text{[cite: 7]}
        \end{center}
    \end{alertblock}
    
    \vspace{0.2cm}
    \textbf{Validity Gate (Para validar la medición empírica)[cite: 7]:}
    \begin{itemize}
        \item ¿Mismo nodo físico de salida ($v_{C2}$) para $A_d$ y $A_c$?[cite: 7]
        \item ¿Misma convención de amplitud en osciloscopio?[cite: 7]
    \end{itemize}
    
    \textbf{Ejemplo de Inconsistencia (No Válido)[cite: 7]:}
    Usar $A_{d,diff} = 80$ (midiendo entre los dos colectores) y $A_{c,se} = 0.04$ (midiendo un solo colector). ¡Las salidas no coinciden![cite: 7]
\end{frame}

\end{document}
